% AUTHOR: Christopher Boerner
% CREATE DATE: 1/8/25
% PURPOSE: LaTex version of resume.
% SPECIAL NOTES: Adapted from Hunter Ellis' resume (https://github.com/hunterwellis/LaTeX-Resume)

\documentclass[9pt, letterpaper]{extarticle}

% Packages:
\usepackage[
    ignoreheadfoot, % set margins without considering header and footer
    top=1.0 cm, % seperation between body and page edge from the top
    bottom=1.0 cm, % seperation between body and page edge from the bottom
    left=1.0 cm, % seperation between body and page edge from the left
    right=1.0 cm, % seperation between body and page edge from the right
    footskip=1.0 cm, % seperation between body and footer
    % showframe % for debugging
]{geometry} % for adjusting page geometry
\usepackage{titlesec} % for customizing section titles
\usepackage{tabularx} % for making tables with fixed width columns
\usepackage{array} % tabularx requires this
\usepackage[dvipsnames]{xcolor} % for coloring text
\definecolor{primaryColor}{RGB}{0, 0, 0} % define primary color
\usepackage{enumitem} % for customizing lists
\usepackage{amsmath} % for math
\usepackage[
    pdftitle={Christopher Boerner's Resume},
    pdfauthor={Christopher Boerner},
    pdfcreator={LaTeX with VSCode},
    colorlinks=true,
    urlcolor=primaryColor
]{hyperref} % for links, metadata and bookmarks
\usepackage[pscoord]{eso-pic} % for floating text on the page
\usepackage{calc} % for calculating lengths
\usepackage{bookmark} % for bookmarks
\usepackage{lastpage} % for getting the total number of pages
\usepackage{changepage} % for one column entries (adjustwidth environment)
\usepackage{paracol} % for two and three column entries
\usepackage{ifthen} % for conditional statements
\usepackage{needspace} % for avoiding page brake right after the section title
\usepackage{iftex} % check if engine is pdflatex, xetex or luatex
\usepackage{datetime}

% Ensure that generate pdf is machine readable/ATS parsable:
\ifPDFTeX
    \input{glyphtounicode}
    \pdfgentounicode=1
    \usepackage[T1]{fontenc}
    \usepackage[utf8]{inputenc}
    \usepackage{lmodern}
\fi

\usepackage{charter}

% Some settings:
\raggedright
\AtBeginEnvironment{adjustwidth}{\partopsep0pt} % remove space before adjustwidth environment
\pagestyle{empty} % no header or footer
\setcounter{secnumdepth}{0} % no section numbering
\setlength{\parindent}{0pt} % no indentation
\setlength{\topskip}{0pt} % no top skip
\setlength{\columnsep}{0.15cm} % set column seperation
\pagenumbering{gobble} % no page numbering

\titleformat{\section}{\needspace{4\baselineskip}\bfseries\Large}{}{0pt}{}[\vspace{1pt}\titlerule]

\titlespacing{\section}{
    % left space:
    -1pt
}{
    % top space:
    0.15 cm
}{
    % bottom space:
    0.15 cm
} % section title spacing

\renewcommand\labelitemi{$\vcenter{\hbox{\small$\bullet$}}$} % custom bullet points
\newenvironment{highlights}{
    \begin{itemize}[
        topsep=0 pt,
        parsep=0 pt,
        partopsep=0pt,
        itemsep=0pt,
        leftmargin=0.25 cm + 10pt
    ]
}{
    \end{itemize}
} % new environment for highlights

\newenvironment{highlightsforbulletentries}{
    \begin{itemize}[
        topsep=0.10 cm,
        parsep=0.10 cm,
        partopsep=0pt,
        itemsep=0pt,
        leftmargin=10pt
    ]
}{
    \end{itemize}
} % new environment for highlights for bullet entries

\newenvironment{onecolentry}{
    \begin{adjustwidth}{
        0 cm + 0.00001 cm
    }{
        0 cm + 0.00001 cm
    }
}{
    \end{adjustwidth}
} % new environment for one column entries

\newenvironment{twocolentry}[2][]{
    \onecolentry
    \def\secondColumn{#2}
    \setcolumnwidth{\fill, 4.5 cm}
    \begin{paracol}{2}
}{
    \switchcolumn \raggedleft \secondColumn
    \end{paracol}
    \endonecolentry
} % new environment for two column entries

\newenvironment{threecolentry}[3][]{
    \onecolentry
    \def\thirdColumn{#3}
    \setcolumnwidth{, \fill, \fill, \fill}
    \begin{paracol}{3}
    {\raggedright #2} \switchcolumn
    \begin{center}  % center the middle column content
}{
    \end{center}
    \switchcolumn \raggedleft \thirdColumn
    \end{paracol}
    \endonecolentry
} % new environment for three column entries

\newenvironment{header}{
    \setlength{\topsep}{0pt}\par\kern\topsep\linespread{1.5}
}{
    \par\kern\topsep
} % new environment for the header

% save the original href command in a new command:
\let\hrefWithoutArrow\href

% ─── DOCUMENT ───────────────────────────────────────────────────────────────
\begin{document}
    \newcommand{\AND}{\unskip
        \cleaders\copy\ANDbox\hskip\wd\ANDbox
        \ignorespaces
    }
    \newsavebox\ANDbox
    \sbox\ANDbox{$|$}

    % HEADER
    \begin{threecolentry}
        {
            Falls Church, Virginia\\
            \hrefWithoutArrow{tel:+1-571-395-5448}{(571) 395-5448}\\
            \hrefWithoutArrow{mailto:boernerc20@gmail.com}{boernerc20@gmail.com}
        }
        {
            Active Top Secret Clearance \\
            Eagle Scout \\
            \hrefWithoutArrow{https://www.boernerc20.me}{boernerc20.me} \\
        }
        {
            \Huge\textbf{{Christopher Boerner}}\\
            \kern 3.0 pt%
            \Large{\textbf{Electrical Engineer}}
        }
    \end{threecolentry}
    \hrule
    \kern 3.0 pt%
    Electrical Engineer at Northrop Grumman supporting the Command \& Data Handling (C\&DH) subsystem on the ESPAStar satellite product line. Computer Engineering Master's graduate with hands-on experience spanning embedded C firmware, custom PCB design, and FPGA/SoC integration. Background includes hardware bring-up and validation, sensor integration, and digital logic implementation.

    % SKILLS SECTION:
    \section{Skills}
    \textbf{Languages \& Software: }C/C++, Python, Verilog, STOL, Bash, MATLAB, Xilinx Vivado/Vitis, STM32CubeIDE, KiCad, LTSpice, Cadence Virtuoso, Altium, Git, Linux\\
    \kern 3.0 pt%
    \textbf{Circuit Design \& EDA: }PCB Layout, Schematic Capture, Signal Conditioning, Power Electronics, FPGA/SoC Integration, UART/SPI/I2C Interfacing\\
    \kern 3.0 pt%
    \textbf{Instrumentation \& Test: }Oscilloscope, DMM, Logic Analyzer, Power Supplies, Scopy, PIL/HIL Validation

    % EDUCATION SECTION:
    \section{Education}
    \begin{twocolentry}{{May 2025}\\\textit{Alexandria, Virginia}}
        \textbf{Master of Engineering in Computer Engineering}\\
        Virginia Tech -- Focused on Computer Systems -- GPA: 3.8\\
        \quad\quad\textit{Advisor: \hrefWithoutArrow{https://www.yangyi.ece.vt.edu/}{{Dr. Cindy Yi (Virginia Tech)}}}\\
    \end{twocolentry}
    \kern 3.0 pt%
    \begin{twocolentry}{{May 2024}\\\textit{Blacksburg, Virginia}}
        \textbf{Bachelor of Science in Computer Engineering}\\
        Virginia Tech -- General Computer Engineering -- GPA: 3.6
    \end{twocolentry}

    % EXPERIENCE SECTION:
    \section{Experience}

    \begin{twocolentry}{{Nov 2025 -- Present} \\\textit{Dulles, Virginia}}
        \textbf{Northrop Grumman | ESPAStar Product Line}\\
        \textbf{\textit{Electrical Engineer}}\\
        \begin{onecolentry}
            \begin{highlights}
                \item Support the Command \& Data Handling (C\&DH) subsystem on the ESPAStar product line -- Northrop Grumman's ESPA-class satellite bus -- and its hosted payloads.
                \item As a Systems Engineer on the NASA Gateway program, analyzed electrical schematics, pinouts, and block diagrams for the Electrical Power System to develop hardware verification test procedures.
                \item Wrote and executed test scripts (Python, NASA STOL) against physical flight computers and flatsat hardware; traced command/telemetry flows through flight software to verify subsystem behavior.
            \end{highlights}
        \end{onecolentry}
    \end{twocolentry}

    \kern 3.0 pt%
    \begin{twocolentry}{{Nov 2024 -- May 2025} \\\textit{Alexandria, Virginia}}
        \textbf{Virginia Tech | FPGA-Accelerated Echo-State Network (Master's Project)}\\
        \textbf{\textit{Graduate Researcher}}\\
        \begin{onecolentry}
            \begin{highlights}
                \item Implemented a floating-point Echo-State Network and online RLS trainer in bare-metal C on a Xilinx Zynq SoC for 2x2 MIMO-OFDM wireless channel prediction.
                \item Integrated FPGA logic with the ARM processing system using Vivado/Vitis; achieved sub-1 ms inference latency with live Ethernet dataset streaming and UART monitoring.
            \end{highlights}
        \end{onecolentry}
    \end{twocolentry}

    \kern 3.0 pt%
    \begin{twocolentry}{{Jun 2024 -- Aug 2024} \\\textit{Grenoble, France}}
        \textbf{Grenoble Electrical Engineering Laboratory | Expe-SmartHouse}\\
        \textbf{\textit{Research Intern}}\\
        \begin{onecolentry}
            \begin{highlights}
                \item Designed photovoltaic-powered smart home mockups with load management and real-time power monitoring.
                \item Programmed Arduino and ESP8266 microcontrollers for MQTT telemetry to a centralized Raspberry Pi energy manager.
            \end{highlights}
        \end{onecolentry}
    \end{twocolentry}

    \kern 3.0 pt%
    \begin{twocolentry}{{Aug 2023 -- May 2024} \\\textit{Blacksburg, Virginia}}
        \textbf{Virginia Tech/NAVAIR | Aircraft Data Acquisition System (Senior Design)}\\
        \textbf{\textit{Team Member}}\\
        \begin{onecolentry}
            \begin{highlights}
                \item Designed custom PCBs with regulated power supplies, Wheatstone bridge signal conditioning, and optocoupler isolation for vibration, temperature, humidity, and acoustic sensor channels.
                \item Programmed STM32 and Artemis MCUs in C to acquire and packetize sensor data; transmitted over RF transceivers to a handheld unit for live telemetry display; validated under MIL-STD-810G.
            \end{highlights}
        \end{onecolentry}
    \end{twocolentry}
    \kern 3.0 pt%

    % PROJECTS SECTION:
    \section{Projects}

    \begin{minipage}[t]{0.48\textwidth}
    \textbf{\large Academic Projects}

    \vspace{4pt}

    \textbf{VLSI Design Project | 12-bit Multiplier in Cadence Virtuoso} \\
    \textit{Nov -- Dec 2023}
    \begin{highlights}
        \item Designed a 12-bit Braun multiplier (schematic/layout) with carry-select adders in Cadence Virtuoso.
        \item Verified via DRC/LVS/PEX; characterized propagation delay, power, and area for ADP optimization.
    \end{highlights}

    \vspace{4pt}

    \textbf{Integrated Design Project | Blood Oxygen Sensor} \\
    \textit{Jan -- May 2022}
    \begin{highlights}
        \item Created a multi-stage amplification and filtration circuit for photodiode sensor signals.
        \item Multiplexed red and infrared channels to calculate blood oxygen saturation.
    \end{highlights}
    \end{minipage}
    \hfill
    \begin{minipage}[t]{0.48\textwidth}
    \textbf{\large Personal Projects}

    \vspace{4pt}

    \textbf{Smart Home Room Node | ESP32-S3, KiCad, Home Assistant} \\
    \textit{Mar -- Present}
    \begin{highlights}
        \item Designing a custom PCB room node integrating voice I/O, temperature/humidity sensing, and an e-ink display with LED ring.
        \item Full hardware cycle from breadboard prototype through KiCad PCB layout, SMD hot-air assembly, and 3D printed enclosure.
    \end{highlights}

    \vspace{4pt}

    \textbf{Distributed Homelab Infrastructure | TrueNAS, Docker, Linux} \\
    \textit{Jan -- Present}
    \begin{highlights}
        \item Self-hosted ZFS NAS, local AI stack (Whisper STT, Ollama LLM), and custom ESP32 room nodes.
        \item Deployed PoE-powered Raspberry Pi nodes for network-wide DNS filtering and Home Assistant smart home automation.
    \end{highlights}
    \end{minipage}

\end{document}
